% SEGMENT OLP-0519-S01
% Part: counterfactuals
% Chapter: introduction
% Section: paradoxes-material

\documentclass[../../../include/open-logic-section]{subfiles}

\begin{document}

\olfileid{cnt}{int}{par}

\olsection{பொருள்சார் நிபந்தனையின் முரண்தோற்றங்கள்}

ஆங்கில ``if \dots then \dots'' கூற்றுகளைக் குறியீடாக்க $!A \lif !B$
ஐப் பயன்படுத்துவதை முதன்முதலில் விமர்சித்தவர்களில் ஒருவர் C.~I. லூயிஸ்.
வைட்ஹெட் மற்றும் ரஸ்ஸலின் \emph{Principia Mathematica} வில் பொருள்சார்
நிபந்தனை பயன்படுத்தப்பட்டதையும், அவர்கள்~$\lif$ ஐ ``implies'' என்று
வாசித்ததையும் லூயிஸ் விமர்சித்தார். $\lif$ என்பது ``பின்விளைகிறது''
என்று பொருள்பட்டால், பொய்யான எந்தக் கூற்று~$p$ உம் $p$ இலிருந்து~$q$
பின்விளைகிறது என்று பின்விளைவிக்கும்; ஏனெனில் $p$ பொய்யெனில் $p \lif
(p \lif q)$ மெய். மேலும், மெய்யான எந்தக் கூற்று~$q$ உம் $p$ இலிருந்து
$q$ பின்விளைகிறது என்று பின்விளைவிக்கும்; ஏனெனில் $q$ மெய்யெனில் $q
\lif (p \lif q)$ மெய். இவை குறித்த லூயிஸின் புகார் சரியானதே.

\emph{உட்கிடைவு}, அதாவது தருக்கமுறைப் பின்விளைவு, ஓர் இணைப்பி அல்ல;
!!{formula}s அல்லது கூற்றுகளுக்கு இடையிலான ஒரு தொடர்பு என்பதைத்
தருக்கவியலாளர்கள் அறிவார்கள். எனவே குழப்பத்தைத் தவிர்க்க $\lif$ ஐ
``பின்விளைகிறது'' என்று வாசிக்கக்கூடாது.\footnote{``$\lif$'' ஐ
  ``பின்விளைகிறது'' என்று வாசிக்கும் பழக்கம் கணிதவியலாளர்களிடமும்
  கணினி அறிவியலாளர்களிடமும் இன்னும் பரவலாக உள்ளது. லூயிஸ் சுட்டிய
  குழப்பங்களைத் தவிர்க்க மெய்யியலாளர்கள் அதை ``only if'' என்று
  வாசிக்க முயல்கின்றனர்.} அப்படி வாசிக்காத வரை, லூயிஸ் எழுப்பிய இந்தக்
குறிப்பிட்ட கவலை எழுவதில்லை: பொய்யாக நிகழும் ஓர் ஆங்கில வாக்கியத்தை~$p$
குறிப்பதாகக் கொண்டாலும், $p$ இலிருந்து $q$ ``பின்விளைவதில்லை.'' $p
\Entails q$ ஆ என்று தீர்மானிக்க \emph{எல்லா} !!{valuation}s ஐயும்
கருத வேண்டும்; நிகழ்வில் பொய்யான ஒரு வாக்கியத்தைக் குறிக்க $p$ ஐப்
பயன்படுத்தினாலும் $p \Entails/ q$.

இருப்பினும், லூயிஸின் பின்வரும் கூற்றைப் போன்ற ``if \dots then \dots''
கூற்றுகளில் ஏதோ விந்தையாக உள்ளது:
\begin{quote}
நிலவு பச்சைப் பாலாடைக்கட்டியால் ஆனது எனில், $2+2=4$.
\end{quote}
பின்வரும் அனுமானங்களிலும் அவ்வாறே உள்ளது:
\begin{quote}
  நிலவு பச்சைப் பாலாடைக்கட்டியால் ஆனதல்ல. எனவே, நிலவு பச்சைப்
  பாலாடைக்கட்டியால் ஆனது எனில், $2+2=4$.

  $2+2 = 4$. எனவே, நிலவு பச்சைப் பாலாடைக்கட்டியால் ஆனது எனில்,
  $2+2=4$.
\end{quote}
ஆயினும், ``if \dots then \dots'' என்பது வெறுமனே~$\lif$ ஆக இருந்தால்,
அந்த வாக்கியம் எந்தச் சிக்கலுமின்றி மெய்யாகவும், அனுமானங்கள் எந்தச்
சிக்கலுமின்றிச் செல்லுபடியாகவும் இருக்கும்.

மற்றோர் எடுத்துக்காட்டு $(!A \lif !B) \lor (!B \lif !A)$ என்ற
மெய்மத்தைப் பற்றியது. செய்தித்தாளிலிருந்து சுட்டுநிலை வாக்கியங்கள்~$S$,
$T$ இரண்டைத் தற்செயலாக எடுத்தால், ``$S$ எனில் $T$, அல்லது $T$ எனில்
$S$'' என்ற வாக்கியம் மெய்யாக இருக்க வேண்டும் என்று இது சுட்டும்.
% SOURCE CORRECTION TA-CF-002: “Another example of concerns” is read as “Another example concerns.”

\end{document}
